% GENERATED FILE. Edit canonical lessons, then run npm run generate:books.
% canonical-source-hash: fnv1a64:ea4a14b7478deeea
% canonical-lessons: RU-C315-ugoshcheniye, RU-C315-ugoshchat, RU-C315-zastolye, RU-C315-chokatsya, RU-C315-shampanskoye

\chapter{Treats and feasts}
\label{ch:ru-treats-and-feasts}
\hlchaptermodality{315}

\begin{chapteropening}
\textbf{By the end of this chapter:} \emph{I can talk about treats, treating guests, a feast, clinking glasses and champagne.}

Complete the last lesson of chapter 315: I can talk about treats, treating guests, a feast, clinking glasses and champagne.
\end{chapteropening}

\section[\texorpdfstring{\ru{угощение} (ugoshchéniye)}{угощение (ugoshchéniye)}]{\ru{угощение} (ugoshchéniye) — treats, refreshments}

\label{lesson:RU-C315-ugoshcheniye}



\begin{quote}
Before the new one: say the Russian for a party popper, then the Russian for a balloon.
\end{quote}

\subsection*{You'll want to know: \ru{угощение}}
\textbf{\ru{угощение}} (\emph{ugoshchéniye}) — \textquotedblleft{}treats, refreshments\textquotedblright{}.

\begin{quote}
\ru{Угощение} \ru{было} \ru{очень} \ru{вкусным}.
\end{quote}

\begin{quote}
\emph{Ugoshchéniye býlo óchen vkúsnym.}
\end{quote}

\begin{quote}
The treats were very tasty.
\end{quote}

Neuter.

\begin{cousinweb}
\emph{\ru{угощение}} is built on \emph{\ru{угощать}} \textquotedblleft{}to treat\textquotedblright{}, from \emph{\ru{гость}} \textquotedblleft{}a guest\textquotedblright{}, inherited from Proto-Slavic \emph{gostĭ}, PIE \emph{g\textsuperscript{h}óstis} (EtymDB).
\end{cousinweb}

\begin{grammarlens}[title={What you've built}]
One more word for telling, explaining and describing at B1.
\end{grammarlens}

\subsection*{Your turn}
\begin{itemize}
\raggedright
  \item \emph{Say it:} \emph{ugoshchéniye}
  \item \emph{Say it:} \emph{ugoshchéniye}, once more
  \item \emph{Say it:} a sentence of your own with \emph{ugoshchéniye}
  \item \emph{Recall:} say the Russian for a party popper, then the Russian for a balloon, then say \emph{ugoshchéniye} again
\end{itemize}

\begin{tcolorbox}[breakable,colback=teal!4,colframe=teal!35!black,title={Before you move on}]
What does \ru{угощение} mean? (\textquotedblleft{}Treats, refreshments\textquotedblright{}.) Say it once more.
\end{tcolorbox}
\section[\texorpdfstring{\ru{угощать} (ugoshchát)}{угощать (ugoshchát)}]{\ru{угощать} (ugoshchát) — to treat (to food or drink)}

\label{lesson:RU-C315-ugoshchat}



\begin{quote}
Before the new one: say the Russian for a balloon, then the Russian for treats.
\end{quote}

\subsection*{You'll want to know: \ru{угощать}}
\textbf{\ru{угощать}} (\emph{ugoshchát}) — \textquotedblleft{}to treat (to food or drink)\textquotedblright{}.

\begin{quote}
\ru{Хозяйка} \ru{угощала} \ru{гостей} \ru{пирогами}.
\end{quote}

\begin{quote}
\emph{Khozyáyka ugoshchála gostéy pirogámi.}
\end{quote}

\begin{quote}
The hostess treated the guests to pies.
\end{quote}

Imperfective; perfective \emph{\ru{угостить}} (\emph{ugostít}), \emph{\ru{угощу}}. The guest is in the accusative, the food in the instrumental.

\begin{cousinweb}
\emph{\ru{угощать}} is built on \emph{\ru{у}-} and \emph{\ru{гость}} \textquotedblleft{}a guest\textquotedblright{}, inherited from Proto-Slavic \emph{gostĭ}, from PIE \emph{g\textsuperscript{h}óstis} (EtymDB): to treat as a guest.
\end{cousinweb}

\begin{grammarlens}[title={What you've built}]
One more word for telling, explaining and describing at B1.
\end{grammarlens}

\subsection*{Your turn}
\begin{itemize}
\raggedright
  \item \emph{Say it:} \emph{ugoshchát}
  \item \emph{Say it:} \emph{ugoshchát}, once more
  \item \emph{Say it:} a sentence of your own with \emph{ugoshchát}
  \item \emph{Recall:} say the Russian for a balloon, then the Russian for treats, then say \emph{ugoshchát} again
\end{itemize}

\begin{tcolorbox}[breakable,colback=teal!4,colframe=teal!35!black,title={Before you move on}]
What does \ru{угощать} mean? (\textquotedblleft{}To treat (to food or drink)\textquotedblright{}.) Say it once more.
\end{tcolorbox}
\section[\texorpdfstring{\ru{застолье} (zastólye)}{застолье (zastólye)}]{\ru{застолье} (zastólye) — a feast, a long festive meal}

\label{lesson:RU-C315-zastolye}



\begin{quote}
Before the new one: say the Russian for treats, then the Russian for to treat.
\end{quote}

\subsection*{You'll want to know: \ru{застолье}}
\textbf{\ru{застолье}} (\emph{zastólye}) — \textquotedblleft{}a feast, a long festive meal\textquotedblright{}.

\begin{quote}
\ru{Застолье} \ru{продолжалось} \ru{до} \ru{утра}.
\end{quote}

\begin{quote}
\emph{Zastólye prodolzhálos do utrá.}
\end{quote}

\begin{quote}
The feast went on until morning.
\end{quote}

Neuter.

\begin{cousinweb}
\emph{\ru{застолье}} is built on \emph{\ru{за}} \textquotedblleft{}at, behind\textquotedblright{} and \emph{\ru{стол}} \textquotedblleft{}table\textquotedblright{}, inherited from Proto-Slavic \emph{stolŭ} (EtymDB): sitting at table.
\end{cousinweb}

\begin{grammarlens}[title={What you've built}]
One more word for telling, explaining and describing at B1.
\end{grammarlens}

\subsection*{Your turn}
\begin{itemize}
\raggedright
  \item \emph{Say it:} \emph{zastólye}
  \item \emph{Say it:} \emph{zastólye}, once more
  \item \emph{Say it:} a sentence of your own with \emph{zastólye}
  \item \emph{Recall:} say the Russian for treats, then the Russian for to treat, then say \emph{zastólye} again
\end{itemize}

\begin{tcolorbox}[breakable,colback=teal!4,colframe=teal!35!black,title={Before you move on}]
What does \ru{застолье} mean? (\textquotedblleft{}A feast, a long festive meal\textquotedblright{}.) Say it once more.
\end{tcolorbox}
\section[\texorpdfstring{\ru{чокаться} (chókatsya)}{чокаться (chókatsya)}]{\ru{чокаться} (chókatsya) — to clink glasses}

\label{lesson:RU-C315-chokatsya}



\begin{quote}
Before the new one: say the Russian for to treat, then the Russian for a feast.
\end{quote}

\subsection*{You'll want to know: \ru{чокаться}}
\textbf{\ru{чокаться}} (\emph{chókatsya}) — \textquotedblleft{}to clink glasses\textquotedblright{}.

\begin{quote}
\ru{Гости} \ru{чокались} \ru{и} \ru{смеялись}.
\end{quote}

\begin{quote}
\emph{Gósti chókalis i smeyális.}
\end{quote}

\begin{quote}
The guests clinked glasses and laughed.
\end{quote}

Imperfective; perfective \emph{\ru{чокнуться}} (\emph{chóknutsya}). Careful: in slang \emph{\ru{чокнутый}} means \textquotedblleft{}crazy\textquotedblright{}.

\begin{cousinweb}
\emph{\ru{чокаться}} is built on an imitative syllable \emph{\ru{чок}} for the ring of glasses; EtymDB lists no parent.
\end{cousinweb}

\begin{grammarlens}[title={What you've built}]
One more word for telling, explaining and describing at B1.
\end{grammarlens}

\subsection*{Your turn}
\begin{itemize}
\raggedright
  \item \emph{Say it:} \emph{chókatsya}
  \item \emph{Say it:} \emph{chókatsya}, once more
  \item \emph{Say it:} a sentence of your own with \emph{chókatsya}
  \item \emph{Recall:} say the Russian for to treat, then the Russian for a feast, then say \emph{chókatsya} again
\end{itemize}

\begin{tcolorbox}[breakable,colback=teal!4,colframe=teal!35!black,title={Before you move on}]
What does \ru{чокаться} mean? (\textquotedblleft{}To clink glasses\textquotedblright{}.) Say it once more.
\end{tcolorbox}
\section[\texorpdfstring{\ru{шампанское} (shampánskoye)}{шампанское (shampánskoye)}]{\ru{шампанское} (shampánskoye) — champagne, sparkling wine}

\label{lesson:RU-C315-shampanskoye}



\begin{quote}
Before the new one: say the Russian for a feast, then the Russian for to clink glasses.
\end{quote}

\subsection*{You'll want to know: \ru{шампанское}}
\textbf{\ru{шампанское}} (\emph{shampánskoye}) — \textquotedblleft{}champagne, sparkling wine\textquotedblright{}.

\begin{quote}
\ru{В} \ru{полночь} \ru{мы} \ru{открыли} \ru{шампанское}.
\end{quote}

\begin{quote}
\emph{V pólnoch my otkrýli shampánskoye.}
\end{quote}

\begin{quote}
At midnight we opened the champagne.
\end{quote}

An adjective serving as a noun, neuter: \emph{\ru{бокал} \ru{шампанского}}.

\begin{cousinweb}
\emph{\ru{шампанское}} is built on the name of the French wine region where it is made, with the Russian adjective ending \emph{-\ru{ское}}; EtymDB lists no parent.
\end{cousinweb}

\begin{grammarlens}[title={What you've built}]
One more word for telling, explaining and describing at B1.
\end{grammarlens}

\subsection*{Your turn}
\begin{itemize}
\raggedright
  \item \emph{Say it:} \emph{shampánskoye}
  \item \emph{Say it:} \emph{shampánskoye}, once more
  \item \emph{Say it:} a sentence of your own with \emph{shampánskoye}
  \item \emph{Recall:} say the Russian for a feast, then the Russian for to clink glasses, then say \emph{shampánskoye} again
\end{itemize}

\begin{tcolorbox}[breakable,colback=teal!4,colframe=teal!35!black,title={Before you move on}]
What does \ru{шампанское} mean? (\textquotedblleft{}Champagne, sparkling wine\textquotedblright{}.) Say it once more.
\end{tcolorbox}
